\section{Kết luận}

Ba mô hình đều học được và chạy inference trên bài toán phát hiện mũ bảo hiểm, nhưng thể hiện ưu/nhược điểm rất khác nhau, và điều này chỉ lộ ra khi tách chỉ số theo từng lớp. Xét chỉ số gộp, Faster R-CNN đứng đầu với mAP@0.5 là 0.650 và mAP@0.5:0.95 là 0.423, cao hơn YOLO (0.582 và 0.389) lẫn RT-DETR (0.555 và 0.369). Tuy nhiên ưu thế đó chủ yếu đến từ lớp \textit{motorbike} --- lớp dễ nhất và ít giá trị nhất với mục tiêu phát hiện vi phạm --- nơi Faster R-CNN đạt AP@0.5 là 0.779 so với 0.518 của YOLO và 0.445 của RT-DETR.

Trên lớp \textit{non-helmet}, tức lớp quyết định của đề tài, thứ hạng đảo lại. RT-DETR đạt AP@0.5 cao nhất là 0.646, YOLO đạt 0.636, còn Faster R-CNN chỉ 0.613. Quan trọng hơn, tại ngưỡng vận hành confidence 0.25, F1 của lớp này là 0.611 với YOLO nhưng chỉ 0.346 với Faster R-CNN và 0.289 với RT-DETR, do precision của YOLO (0.613) vượt xa hai mô hình còn lại (0.238 và 0.180). Kết hợp với tốc độ 36.4 FPS so với 12.4 và 2.8, YOLO là lựa chọn triển khai hợp lý nhất --- không phải vì dẫn đầu mọi chỉ số gộp, mà vì cho chất lượng tốt nhất trên đúng lớp mà hệ thống cần và đáp ứng được yêu cầu thời gian thực.

Kết quả cho thấy khó khăn lớn nhất của bài toán nằm ở dữ liệu: lớp \textit{non-helmet} ít hơn, vật thể cần nhận diện nhỏ và phân phối giữa validation/test chưa hoàn toàn đồng nhất. Vì vậy, ngoài việc điều chỉnh mô hình, hướng cải thiện quan trọng là cân bằng lại dữ liệu, kiểm tra chất lượng pseudo-labels, tăng thêm ảnh chứa vi phạm không đội mũ và đánh giá lại trên split được chia theo nguồn ảnh một cách nhất quán.

Nếu mục tiêu là hệ thống có thể chạy nhanh và cho kết quả ổn định, YOLO là lựa chọn hợp lý nhất trong ba mô hình hiện tại. RT-DETR có tiềm năng tiếp tục cải thiện nếu được tối ưu tốc độ, train với batch lớn hơn, ảnh kích thước cao hơn hoặc dữ liệu cân bằng hơn. Faster R-CNN nên được xem là baseline tham chiếu và cần kiểm soát overfitting nếu tiếp tục phát triển.
